% !TEX root = /home/fred-olav/afgv/src/preamble.tex
\input preamble.tex
\huge
\centerline{\bf Juleprøve 02 - Automatiseringsystemer}  \bigskip
\normalsize


\vskip 3cm 
Kompetansemål:
\begin{itemize}[noitemsep]

	\item Prøven dekker alle kompetansemål i automatiseirngsfaget fra VG1 til VG3 auto.  
\end{itemize}

Alle ark som leveres inn skal ha elevens navn. \\ 

Oppgave 1-6 leveres på papir. NB. riv av forsiden slik at du har denne informasjonen etter innlevering av oppgaver på ark. \\

Oppgave 7-9 skal gjøres på PC og det skal levers som en PDF fil . Denne sendes på mail til fred-olav.mosdal@skole.rogfk.no med emne:
\vskip 5pt 
Juleprøve 

\bigskip 
\hrule
\vfil \eject
\hrule

\bigskip 
\hrule
\vfil \eject
\textbf{Oppgave 1 - Elektroteknikk}
\vskip 5pt
\hrule
\vskip 5pt
% Elektroteknikk
\textbf{a)} Regn ut resistansen til strekklappen i denne målebro kretsen, gitt voltmeterets indikasjon på 1.76228mV (A er positiv og B er negativ).

$$\includegraphics[width=15.5cm]{i04867x01.eps}$$

Gitt at dette er et PT100 element fra europa, hvilken temperatur måler elementet?

\vskip 1cm 
\begin{tikzpicture}
	\draw[step=0.5cm,gray!20,very thin]  grid (16,8) ;
\end{tikzpicture}

\begin{tikzpicture}
	\draw[step=0.5cm,gray!20,very thin]  grid (16,4) ;
\end{tikzpicture}

\vfil \eject
\textbf{b)} Regn ut spenningsfallet i denne 4-20 mA strømsløyfekretsen for strømmene i tabellen:
$$\includegraphics[width=15.5cm]{i00129x01.eps}$$

% No blank lines allowed between lines of an \halign structure!
% I use comments (%) instead, so that TeX doesn't choke.

$$\vbox{\offinterlineskip
\halign{\strut
\vrule \quad\hfil # \ \hfil & 
\vrule \quad\hfil # \ \hfil & 
\vrule \quad\hfil # \ \hfil & 
\vrule \quad\hfil # \ \hfil \vrule \cr
\noalign{\hrule}
%
% First row
Percent of range & Transmitter current & $V_{AB}$ &  $V_{EF}$ \cr
%
\noalign{\hrule}
%
% Another row
0 \% & 4 mA &  &  \cr
%
\noalign{\hrule}
%
% Another row
%
%
% Another row
50 \% & 12 mA &  &  \cr
%
\noalign{\hrule}
%
% Another row
%
%
% Another row
100 \% & 20 mA &  &  \cr
%
\noalign{\hrule}
} % End of \halign 
}$$ % End of \vbox

\vskip 5pt 
\begin{tikzpicture}
	\draw[step=0.5cm,gray!20,very thin]  grid (16,10) ;
\end{tikzpicture}
\filbreak


\vfil \eject
%\vfil
\textbf{Oppgave 2 - Koblinger}
% Koblinger
\vskip 5pt
\hrule

\vskip 5pt
\textbf{a)} Tegn inn ledninger slik at de to parallelle motstandene kommer i serie med den enkle motstanden. Koble så kretsen til batteriet.

$$\includegraphics[width=15cm]{i04814x03.eps}$$
\vskip 1cm
\vfil \eject
\textbf{b)} Tegn inn ledninger sånn at NO brytere vil aktivere releet når den trykkes inn. Releet skal igjen gi strøm til lampen slik at den lyser.

$$\includegraphics[width=15cm]{i04814x04.eps}$$

\vskip 1cm
\filbreak
\textbf{c)} Tegn inn legninger som viser hvordan den kontaktløse radar nivåmåleren skal kobles for og sende prosessveriden til inngang 1 på regulatoren. Du kan anta at strømsløyfen på nivåmåleren er av 4 leder type.
$$\includegraphics[width=15cm]{i04814x05.eps}$$

\vfil \eject
\textbf{Oppgave 3 - Instrumentering}
\vskip 5pt
\hrule
\vskip 5pt

\textbf{a )} I denne kretsen gir en DP-celle signal med 4-20mA til: en lokal trykkindikator og til en fjern trykk indikator. Din oppgave er å finne en måte å koble loop-kalibratoren slik at bare den fjerne trykkindikatoren viser et trykk på 300 bar. Din Kobling skal ikke forstyre den lokale visning av trykk. 


$$\includegraphics[width=15.5cm]{i00627x01.eps}$$

Tegn inn din løsning der du viser loop-kalibratorens testlegninger, strømmen den må sende og moduset den må stå i({\it source}, {\it read}, or {\it simulate}). Anta et måleområde på $-10$ Bar til +440 Bar.  


\vfil \eject

\textbf{b) }En SMART DP-celle er tatt ut av drift og tatt med for benkkalibrering. En automatikker kobler til et presisjons trykkmanometer og en luftkilde til High inngangen på DP-cellen, mens han måler strømutgangen med et multimeter. 

$$\includegraphics[width=15.5cm]{i04815x01.eps}$$

Regn ut avviket i \% av måleområde for {\it sensor trim} og avviket i  \% av måleområde for {\it utgangstrim}.  Forklar hvorfor en må ha en HART kommunikator for å kunne regne disse avvikene sepparat. 

\vskip 1cm 
\begin{tikzpicture}
	\draw[step=0.5cm,gray!20,very thin]  grid (16,5) ;
\end{tikzpicture}



\vfil \eject
\vskip 1cm 
\begin{tikzpicture}
	\draw[step=0.5cm,gray!20,very thin]  grid (16,25) ;
\end{tikzpicture}
\vfil \eject
\textbf{Oppgave 4 - PLS}
% PLS live view og koblinger o.l.
\vskip 5pt
\hrule
\vskip 5pt
\textbf{a)} Du skal tegne inn kobling til NO/NC på bryterne slik at "live" view bilde stemmer. 

\vskip 5pt 
Trykket på byrter A er 14 PSI og på bryter B er det 70  PSI. 

$$\includegraphics[width=15.5cm]{i02882x01.eps}$$

\vskip 1cm
\vfil \eject
\textbf{b) }To endebrytere skal tilkobles h.h.v. \texttt{DI2} og \texttt{DI5} på en Wago 750-436 DI inngangsmodul. Tegn de nødvendige koblingene. Det internekoblingsskjemaet for (\texttt{DI1}) vises som en referanse for alle inngangene.
Tegn de nødvendige koblingene som er nødvendig for tilkobling av endebryterene.  


$$\includegraphics[width=15.5cm]{i04807x01.eps}$$

\vfil \eject
\textbf{Oppgave 5 - Målesystem for nivå}
%Måleteknikk 
\vskip 5pt
\hrule
\vskip 10pt 

Se tegningsnummer 00-IH-343 i forberedelsesdelen til prøven (Sløyfe 343). Hvilket nåleprinsipp er brukt for å måle nivået i denne tanken. Forklar også virkemåten til prinsippet. 

\vskip 10pt 
\begin{tikzpicture}
	\draw[step=0.5cm,gray!20,very thin]  grid (16,23) ;
\end{tikzpicture}


\vfil \eject
\textbf{Oppgave 6 - Målesystem for temperatur}
\vskip 5pt
\hrule
\vskip 10pt

Se tegningsnummer 00-IH-302 i forberedelsesdelen (sløyfe 302). Hvilken temperaur måle metode kunne vært aktuelt å bruke her? Forklar også virkemåten til et målesysem med denne målemetoden. 

\vskip 10pt
\begin{tikzpicture}
	\draw[step=0.5cm,gray!20,very thin]  grid (16,22) ;
\end{tikzpicture}
\vfil \eject
\textbf{Oppgave 7 - Tegneoppgave}
%Tegneoppgave
\vskip 5pt
\hrule
\vskip 5pt
Tegn styrestrømsskjemaet fra tegningen på neste side i PC|Schematic

\vfil \eject
\includepdf[pages=-,angle=90]{../output/Jul21/Skjema.pdf}
\vfil \eject
\vskip 5pt

\textbf{Oppgave 8 - RIO oppgave}
%RIO kommunikasjon oppgave
\vskip 5pt
\hrule

\vskip 5pt
\textbf{a) }Du skal skal sette opp kommunikasjon fra PLS som styrer anlegget rundt H-340 til en sentral HMI skjerm. skjermen kommunikserer på Modbus TCP og har adresse 192.168.0.32. Sett opp en blokkskjematisk oversikt over hvordan komponentene kobles sammen. 
\vskip 10pt 
\textbf{b) }Sett opp et codesys program som viser oppsett av modbus master og modbus slaven. Vær nøye med å skille de to systemene med ulike POU-er.
\small

\begin{center}
\begin{tabular}{ | m{2cm} | m{2.5cm}| m{11cm} |} 
\hline
\multicolumn{3}{|c|}{Modbus data table} \\
\hline
\multicolumn{3}{|c|}{Input registers}\\
\hline
Register No.&\multicolumn{2}{c|}{Innhold} \\
\hline
0000H&TIC342 PV&Temperatur i sløyfe 342\\
\hline
0001H&TIC342 MV&Pådrag i sløyfe 342\\
\hline
	0002H&PIC341 PV &Trykk i sløyfe 341\\
\hline
	0003H&PIC341 MV& Pådrag i sløyfe341\\
\hline
\multicolumn{3}{|c|}{Holding registers}\\
\hline
Register No.&\multicolumn{2}{c|}{Innhold} \\
\hline
	0000H&TIC342 SP& Settpunkt i sløyfe 342\\
\hline
	0001H&PIC341 SP&Settpuknt i sløyfe 341 (0-100\%)\\
\hline
	\multirow{2}{*}{0002H}&bit0&Startsignal P-342\\
	\cline{2-3}
&bit1&Startsignal I321\\
\hline
\end{tabular}
\end{center}
\normalsize


\filbreak
\vfil \eject
\textbf{Oppgave 9}
\vskip 5pt
\hrule
\vskip 5pt
Du får i oppgave å erstatte TIC-302 med ny løsning som vil si bruk av RIO (Remote In Out). 
\vskip 5pt 
Vis hvordan du ville planlagt, gjennomført og dokumentert jobben. 


\vfil \eject
\includepdf[pages=-,angle=90]{../eq/afgvformler.pdf}
\end {document}
